\naslov{Nelinearne enačbe}

Iščemo rešitve enačbe $f(x) = 0$ za nek $f : \R \to \R$.
Pri tem lahko pridemo do več različnih situacij glede obstoja in enoličnosti
rešitve; možnost je, da rešitev obstaja in je ena sama, da je rešitev več, ampak
končno mnogo, da jih je neskončno mnogo, ali pa da rešitve sploh ni.

Naj bo $\alpha$ ničla za zvezno odvedljivo $f$.
Ničla je enostavna natanko tedaj, ko je $f'(\alpha) \ne 0$.
Če ni enostavna, je $m$-kratna natanko tedaj, ko je prvi neničelni odvod v
$\alpha$ reda $m$.
V primeru enostavne ničle lokalno obstaja inverzna funkcija, da je $\alpha =
f^{-1}(0)$.
Absolutna občutljivost problema je tedaj enaka $\abs{(f^{-1})'(0)} =
\inv{\abs{f'(\alpha)}}$.
Če je $\alpha$ dvojna ničla, uporabimo Taylorjev približek druge stopnje
\[
  f(x) \approx
  \underbrace{f(\alpha)}_{=0} + \underbrace{f'(\alpha)}_{=0} (x-\alpha) +
  \inv{2} f''(\alpha) {(x-\alpha)}^2,
\]
torej za $\abs{f(x)} \le \varepsilon$ velja
\[
  \abs{x - \alpha} \le \sqrt{ \frac{2 \varepsilon}{\abs{f''(\alpha)}} }.
\]
Višje kot so ničle, bolj občutljiv je problem iskanja.
Za večkratno ničlo v splošnem velja
\[
  \abs{x - \alpha} \le {\left( \frac{m! \varepsilon}{\abs{f^{(m)}(\alpha)}}
	\right)}^{1/m}
\]

\vprasanje{Analiziraj problem iskanja ničel.}

\podnaslov{Bisekcija}

Pri implementaciji bisekcije si hranimo zaporedja ${(a_n)}_n$ levih mej,
${(b_n)}_n$ desnih mej, ${(c_n)}_n$ sredinskih približkov, in ${(e_n)}_n$
polovičnih velikosti intervala.
Nove člene izračunamo po predpisih
\begin{gather*}
  e_{n+1} = e_n / 2, \\
  a_n = a_{n-1}\ \text{ali}\ c_{n-1}, \\
  b_n = b_{n-1}\ \text{ali}\ c_{n-1}, \\
  c_n = a_n + e_n. \\
\end{gather*}
Pri tem zmanjšamo število računskih operacij, se izognemo problemom glede možnih
nepredvidenih zaokrožitev, skokov izven območja ali računskih napak.
Bisekcija nam lahko poišče liho ničlo, ne pa sode, prav tako lahko poišče lih
pol (ne pa sodega).

\vprasanje{Opiši delovanje bisekcije.}

\podnaslov{Navadna iteracija}

Pri navadni iteraciji iskanje rešitve enačbe $f(x) = 0$ prevedemo na iskanje
rešitve enačbe $x = g(x)$ za ustrezno izbrano funkcijo $g$.
Splošna primerna izbira je recimo $g(x) = x - f(x)$, ali pa $g(x) = x - h(x)
f(x)$ za neničelno funkcijo $h$.
Da postopek $x_{r+1} = g(x_r)$ konvergira, mora biti $g$ v okolici $\alpha$
skrčitev.

\begin{izrek}
  Naj bo $\alpha = g(\alpha)$ in naj $g$ na intervalu $I = [\alpha - \delta,
  \alpha + \delta]$ za nek $\delta > 0$ zadošča Lipschitzovem pogoju $\abs{g(x)
	- g(y)} \le m \abs{x - y}$ za nek $m \in [0, 1)$, in poljubna $x, y \in I$.
  Potem za vsak $x_0 \in I$ zaporedje $x_{r+1} = g(x_r)$ konvergira k $\alpha$,
  in velja
  \begin{itemize}
  \item $\abs{x_r - \alpha} \le m^r \abs{x_0 - \alpha}$,
  \item $\abs{x_{r+1} - \alpha} \le \frac{m}{1-m} \abs{x_{r+1} - x_r}$.
  \end{itemize}
\end{izrek}

\begin{proof}
  Dokažimo prvo, da zaporedje ne zapusti intervala $I$;
  če velja $x_r \in I$, je $\abs{x_r - \alpha} \le \delta$, torej
  \[
	\abs{x_{r+1} - \alpha} = \abs{g(x_r) - g(\alpha)}
	\le m \abs{x_r - \alpha} < \abs{x_r - \alpha}
	\le \delta,
  \]
  torej je tudi $x_{r+1} \in I$.
  S ponavljanjem tega postopka tudi dokažemo prvo točko, za drugo točko pa
  ocenimo
  \begin{align*}
	\abs{x_{r+1} - \alpha}
	&= \abs{x_{r+1} - x_{r+2} + x_{r+2} - x_{r+3} + x_{r+3} - \ldots - \alpha} \\
	&\le \abs{x_{r+1} - x_{r+2}} + \abs{x_{r+2} - x_{r+3}} + \ldots \\
	&\le (m + m^2 + m^3 + \ldots) \abs{x_{r+1} - x_r} \\
	&= \frac{m}{1-m} \abs{x_{r+1} - x_r}.
  \end{align*}
\end{proof}

\begin{posledica}
  Če je $\alpha = g(\alpha)$, če je $g$ zvezno odvedljiva in če velja
  $\abs{g'(\alpha)} < 1$, potem obstaja nek $\delta > 0$, da za vsak $x_0$,
  $\abs{x_0 - \alpha} < \delta$, zaporedje $x_{r+1} = g(x_r)$ konvergira k
  $\alpha$.
\end{posledica}

\vprasanje{Povej in dokaži izrek o navadni iteraciji.}

\begin{definicija}
  Naj bo $\lim x_r = \alpha$.
  Pravimo, da ${(x_r)}_r$ \pojem{konvergira k $\alpha$ z redom $p$}, če velja
  \[
	\lim_{r \to \infty} \frac{\abs{x_{r+1} - \alpha}}{\abs{x_r - \alpha}^p} = C
  \]
  za neko konstanto $C > 0$.
\end{definicija}

\begin{izrek}
  Naj bo $\alpha = g(\alpha)$ za $p$-krat zvezno odvedljivo funkcijo $g$ in naj
  velja $g'(\alpha) = \ldots = g^{(p-1)}(\alpha) = 0$ ter $g^{(p)}(\alpha) \ne
  0$.
  Tedaj v bližini $\alpha$ zaporedje $x_{r+1} = g(x_r)$ konvergira k $\alpha$ z
  redom $p$.
\end{izrek}

\begin{proof}
  Izraz $x_{r+1} = g(x_r)$ v okolici $\alpha$ razvijemo v Taylorjevo vrsto:
  \[
	x_{r+1} = g(\alpha) + g'(\alpha) (x_r - \alpha) + \ldots +
	\frac{g^{(p-1)}(\alpha)}{(p-1)!}(x_r - \alpha)^{p-1}
	+ \frac{g^{(p)}(\xi)}{p!} (x_r - \alpha)^p
  \]
  za nek $\xi$ v bližini $\alpha$.
  Sledi
  \[
	\frac{x_{r+1}-\alpha}{(x_r - \alpha)^p} = \frac{g^{(p)}(\xi)}{p!}.
  \]
\end{proof}

\vprasanje{Kaj je red konvergence zaporedja? Kako ga poiščeš z odvodom?}

\podnaslov{Tangentna metoda}

Če vzamemo $\alpha = x_r + \Delta x_r$, in razvijemo dva člena Taylorjeve vrste
\[
  0 = f(x_r + \Delta x_r) = f(x_r) + f'(x_r) \Delta x_r + \pol f''(\xi_r) \Delta x_r^2,
\]
ter nato zanemarimo zadnji člen, dobimo
\[
  \Delta x_r = \frac{- f(x_r)}{f'(x_r)},
\]
s čimer smo izpeljali tangentno metodo, ki je pravzaprav poseben primer naravne
iteracije.
Če je $\alpha$ enostavna ničla, lahko izračunamo, da ima $g(x) = x - f(x) /
f'(x)$ ničelni odvod v $\alpha$ (ob predpostavki $f \in \zvodvi{2}$), in je
torej konvergenca vsaj kvadratična.
V primeru $m$-kratne ničle za $m \ge 2$ pa po dolgopisni izpeljavi dobimo, da je
konvergenca zagotovljena, a linearna.
Za dvakrat zvezno odvedljive funkcije je vsaka ničla $f$ torej privlačna negibna
točka.

\vprasanje{Izpelji tangentno metodo in pokaži, kakšen red konvergence ima.}

\begin{izrek}
  Naj bo $f$ na $I = [0, \infty)$ dvakrat zvezno odvedljiva, strogo naraščajoča,
  konveksna in naj ima na $I$ ničlo $\alpha$.
  Potem za vsak $x_0 \in I$ tangenta metoda konvergira k $\alpha$.
\end{izrek}

\begin{proof}
  Velja $f'(x) > 0$ in $f''(x) > 0$.
  Z oceno
  \[
	x_1 - \alpha = \frac{f''(\xi)}{2f'(x_0)} (x_0 - \alpha)^2 \ge 0
  \]
  pokažemo, da je ne glede na $x_0$ točka $x_1$ vedno desno od $\alpha$.
  Dokažimo še, da je $x_{r+1}$ nujno med $\alpha$ in $x_r$:
  \[
	x_{r+1} = x_r - \frac{f(x_r)}{f'(x_r)} < x_r,
  \]
  vedno pa velja $x_r > \alpha$.
  To je torej strogo padajoče omejeno zaporedje, ki mora nekam konvergirati; to
  bo seveda $\alpha$.
\end{proof}

\vprasanje{Za kakšne funkcije lahko globalno zagotoviš konvergenco tangentne
  metode? Dokaži.}

\podnaslov{Sekantna metoda}

Če je izračun odvoda zahteven, ga lahko aproksimiramo z diferenčnim kvocientom.
Namesto tangente na $f$ v točki $x_r$ tako uporabimo sekanto skozi točki $(x_r,
f(x_r))$ in $(x_{r-1}, f(x_{r-1}))$.
Dobljena metoda tehnično ni navadna iteracija, ker uporablja zadnja dva
približka, vendar se obnaša sorodno.
Naslednji približek izračunamo s predpisom
\[
  x_{r+1} = x_r - \frac{f(x_r) (x_r - x_{r-1})}{f(x_r) - f(x_{r-1})}.
\]
Analiza sekantne metode je težja kot analiza metod navadne iteracije.
Izkaže se, da velja
\[
  \abs{e_{r+1}} \approx c \abs{e_r} \abs{e_{r-1}}
\]
za neko konstanto $c$, ki je pravzaprav enaka
\[
  c = \frac{\abs{f''(\alpha)}}{2 \abs{f'(\alpha)}}.
\]

Označimo red sekantne metode s $p$.
Obstaja konstanta $D > 0$, da je $\abs{e_{r+1}} \approx D \abs{e_r}^p$,
torej
\[
  \abs{e_{r+1}} \approx C D \abs{e_{r-1}}^{p+1} = D^{p+1} \abs{e_{r-1}}^{p^2},
\]
iz česar sledi $p^2 = p+1$ oziroma $p = \phi$, torej je konvergenca
superlinearna.

\vprasanje{Razloži sekantno metodo in izpelji njen red konvergence.}

\podnaslov{Ostale metode}

Pri \emph{Mullerjevi metodi} uporabimo tri približke $x_r$, $x_{r-1}$ in
$x_{r-2}$, ter skozi točke $(x_r, f(x_r))$, $(x_{r-1}, f(x_{r-1}))$, $(x_{r-2},
f(x_{r-2}))$ potegnemo polinom stopnje 2.
Za naslednji približek vzamemo tisto izmed dveh ničel polinoma, ki je bližnja
$x_r$.
Ena od prednosti te metode je, da lahko išče kompleksne ničle tudi z realnimi
začetnimi približki.
Izkaže se, da je red konvergence približno $1.84$.

\vprasanje{Razloži Mullerjevo metodo.}

Če zamenjamo vlogi $x$ in $y$, in najdemo polinom $p(y)$, ki poteka skozi točke
$(f(x_r), x_r)$, $(f(x_{r-1}), x_{r-1})$, $(f(x_{r-2}), x_{r-2})$, dobimo približek
za $f^{-1}$, in lahko za naslednji približek vzamemo $x_{r+1} = p(0)$.
Metodo imenujemo \emph{inverzna interpolacija}, ima pa isti red konvergence kot
Mullerjeva metoda.

\vprasanje{Razloži metodo inverzne interpolacije.}

\podnaslov{Ničle polinomov}

Ničle polinomov lahko iščemo na več načinov:
\begin{itemize}
\item Poiščeš eno ničlo in reduciraš polinom.
\item Računaš vse ničle hkrati.
\item Prevedeš problem na problem iskanja lastnih vrednosti.
\end{itemize}
Za redukcijo na problem lastnih vrednosti uporabljamo \emph{spremljevalno
  matriko} polinoma $p(x) = a_0 x^n + \ldots + a_n$
\[
  C_p =
  \begin{bmatrix}
	0 & 1 &  & \\
	& 0 & 1 & \\
	&  & \ddots & \ddots &  \\
	& &  & 0 & 1 \\
	- \frac{a_n}{a_0} & - \frac{a_{n-1}}{a_0} & \ldots & -\frac{a_2}{a_0} & -\frac{a_1}{a_0}
  \end{bmatrix}
\]

\vprasanje{Kako izgleda spremljevalna matrika polinoma?}

Ena od metod, ki računa vse ničle hkrati, je \emph{Laguerrova metoda}.
Za polinom $p(x) = a_0 x^n + \ldots + a_n$ z ničlami $\alpha_1, \ldots,
\alpha_n$ definiramo
\begin{gather*}
  S_1(x) = \sum_{i=1}^n \inv{x-\alpha_i} = \frac{p'(x)}{p(x)}, \\
  S_2(x) = \sum_{i=1}^n \inv{{(x-\alpha_i)}^2} = -S_1'(x) = \frac{(p'(x))^2 -
	p(x) p''(x)}{p^2(x)}, \\
  a(x) = \inv{x - \alpha_n}, \\
  b(x) = \inv{n-1} \sum_{i=1}^{n-1} \inv{x - \alpha_i},
\end{gather*}
da velja $S_1(x) = a(x) + (n-1) b(x)$.
Tedaj za
\begin{gather*}
  d_i(x) = \inv{x-\alpha_i} - b(x), \\
  d(x) = \sum_{i=1}^{n-1} d_i^2(x)
\end{gather*}
dobimo $S_2 = a^2 + (n-1) b^2 + d$, ker je $\sum_i d_i = 0$.
Dobili smo sistem enačb v spremenljivkah $a, b$, ki ga lahko rešimo in dobimo
\[
  a_{1, 2} = \inv{n} \left( S_1 \pm \sqrt{(n-1) (n S_2 - S_1^2 - n d)} \right).
\]
Če $x$ obravnavamo kot približek za ničlo $\alpha_n$, bo člen $nd$ v bližini
$\alpha_n$ majhen, zato ga zanemarimo.
Iz tega izrazimo
\[
  \alpha_n = x - \frac{n}{S_1 \pm \sqrt{ (n-1)(n S_2 - S_1^2) }}.
\]
Laguerrova metoda nam torej da postopek za izračun približka ničle
\[
  x_{r+1} = x_r -
  \frac{n p(x_r)}{
	p'(x_r) \pm
	\sqrt{ (n-1) \left[ (n-1) (p'(x_r))^2 - n p(x_r) p''(x_r) \right] }
  }.
\]
Za odločitev, kateri predznak pripišemo korenu v imenovalcu, imamo tri možnosti:
\begin{itemize}
\item Vedno izberemo plus,
\item Vedno izberemo minus,
\item \emph{Stabilna varianta}: izbereš tistega, ki ti v imenovalcu da večjo
  absolutno vrednost.
\end{itemize}
V prvih dveh primerih fiksno iščemo v eni smeri od začetnega približka, v
tretjem pa to ni zagotovljeno.

\begin{izrek}
  Če ima polinom $p$ same realne ničle, potem za vsak začetni približek $x_0$
  stabilna verzija Laguerrove metode konvergira proti najbližji desni oz. levi
  ničli, pri čemer si mislimo, da sta kraka realne osi pri $+\infty$ in
  $-\infty$ združena.
  Konvergenca v bližini enostavne ničle je kubična.
\end{izrek}

Če ima polinom kompleksne ničle, metoda konvergira za skoraj vse začetne
približke.

\vprasanje{Izpelji Laguerrovo metodo in razloži, kako delujejo vse možnosti za
  izbiro naslednjega približka.}

\podnaslov{Durand-Kernerjeva metoda}

Izberimo približke $x_1, \ldots, x_n$ za ničle polinoma $p(x)$ z vodilnih
koeficientom $1$.
Iščemo popravke $\Delta x_1, \ldots, \Delta x_n$, da bodo $x_i + \Delta x_i$
točne ničle.
Velja
\begin{align*}
  p(x)
  &= (x - (x_1 + \Delta x_1)) (x - (x_2 + \Delta x_2)) \ldots (x - (x_n
  + \Delta x_n)) \\
  &= \prod_{i=1}^n (x-x_i) - \sum_{j=1}^n \Delta x_j \prod{i \ne j} (x - x_i)
  + \ldots,
\end{align*}
člene drugega in večjega reda pa zanemarimo (torej vse, kar je v tropičju).
Če s $q(x)$ označimo nezanemarjen del, velja
\[
  q(x_l) = - \Delta x_l \prod_{i \ne l} (x_l - x_i),
\]
iz česar lahko izračunamo $\Delta x_l$.

\vprasanje{Razloži Durand-Kernerjevo metodo.}